%%%%%%%%%%  scan idx 82  |  folio 69  %%%%%%%%%%
\section*{\S 6. Le théorème de Riemann--Roch}

Son énoncé est le suivant~: \emph{soit \(f\) un morphisme} propre \emph{d'une variété quasi-projective non singulière \(Y\) dans une variété \(X\) de même espèce. Alors, on a~:}
\begin{equation}\tag{2.39}
 f_*\bigl(\ch_Y(x)\,\mathcal T(Y)\bigr)
 =
 \ch_X\bigl(f_*(x)\bigr)\,\mathcal T(X).
\end{equation}
(Les deux membres sont dans \(A(X)\otimes\mathbb Q\) et \(\mathcal T(X)=\mathcal T(C(T_X))\).)

Il permet de calculer, aux éléments de torsion près, les classes de Chern de \(f_!(x)\) connaissant les classes de Chern de \(x\), celles de \(T_X\) et \(T_Y\) et l'homomorphisme \(f_*\). Ou encore, il permet de déterminer l'homomorphisme
\(f_!\) de \(\underline K(Y)\otimes\mathbb Q\) dans \(\underline K(X)\otimes\mathbb Q\), lorsqu'on identifie au moyen des homomorphismes \(\ch_Y\) et \(\ch_X\) ces groupes respectivement à \(A(Y)\otimes\mathbb Q\) et \(A(X)\otimes\mathbb Q\) (si la formule (2.35) est vraie~!), à l'aide de \(f_*\) et des classes de Chern de \(T_X\) et de \(T_Y\).

Je ne sais démontrer la formule (2.39) qu'en caractéristique \(0\)\footnote{Voir plus bas une démonstration valable en toutes caractéristiques, basée sur un principe différent (p.\,51).}, et comme égalité dans \(G(\underline K(X))\otimes\mathbb Q\). Pour ceci \(Y\) étant localement fermé dans un espace projectif \(P\), on décompose \(f\) en le produit d'un morphisme d'injection de \(Y\) dans \(P\times X\), et d'un morphisme de projection de \(P\times X\) sur \(X\), et l'on démontre (2.39) pour chacun de ces deux morphismes..

Pour un morphisme d'injection, on voit immédiatement que les formules (2.38) et (2.39) sont équivalentes (l'hypothèse de caractéristique \(0\) est intervenue uniquement par l'intermédiaire du théorème 2.8 qui permet d'exprimer les \(\lambda^i(\gamma(F))\), pour un faisceau algébrique cohérent \(F\), par des invariants cohomologiques \emph{locaux} de \(F\)).

Pour un morphisme de projection \(f:P\times X\longrightarrow X\), on peut, grâce à la proposition 2.13, supposer que \(x\) est de la forme \(y\otimes z\), avec \(y\in\underline K(P)\) et 
